\subsection{%Essential Euler-Lagrange estimate:
Proof of Lemma~\ref{la:lhs}}

The main idea is to use $\phi$ as the derivative of a test function for the 
Euler-Lagrange equation. Of course this is not possible generally, but with some
precaution we can actually do it up to a benign error term:
For $\phi \in C^\infty_0 ((4/10,6/10))$ we set
\begin{equation*}
   h(u):= \int_0^u \phi(v) dv - a u 
\end{equation*}
where $a=\int_0^1 \phi(v) dv$  is chosen such that $ h(0)= h(1)$. 
Hence, if we set
\begin{align*}
h_\pi(k+t) := h(t), \quad \forall t\in [0,1) , k \in \mathbb Z, \\
\phi_\pi(k+t) := \phi(t), \quad \forall t\in [0,1) , k \in \mathbb Z,
\end{align*}
$h_\pi$ is a smooth one periodic function satisfying
\begin{equation*}
    h_\pi'(u) =  \phi_\pi  - a.
\end{equation*}
and, as we assume in Lemma~\ref{la:lhs} that $\left\langle\phi,\gamma'\right\rangle \equiv 0$, also
\begin{equation*}
 \left\langle \phi_\pi , \gamma'\right \rangle \equiv 0.
\end{equation*}

Since $\gamma$ is a stationary point of the M\"obius energy, testing the 
equation \eqref{eq:knotPDE} with $h_\pi$, (recall \eqref{eq:defT1}, \eqref{eq:defT2},\eqref{eq:defQeps}) 
\begin{equation}\label{eq:ourELeqn}
 Q(\gamma,h_\pi) := \lim_{\varepsilon \to 0} Q_\varepsilon(\gamma,h_\pi)
= T_1(\gamma,h_\pi) + T_2(\gamma,h_\pi).
\end{equation}

As $\phi_\pi$ is perpendicular to $\gamma'$ we can estimate the term $T_1(\gamma,h_\pi)$
by
\begin{align}
 T_1(\gamma,h_\pi) &= \int_{-1}^1 \int_{-1/2}^{1/2}
   \left\langle \gamma'(u),\phi_\pi(u) - a \right\rangle 
  \left(\frac 1 {|\gamma(u+w)- \gamma(u)|^2} - \frac 1 {w^2} \right)\nonumber\\
  &\leq |a|\ E^{(2)}(\gamma) \leq C \|\phi\|_{L^1} \leq C \|\phi\|_{L^2}.\label{eq:toRT1}
\end{align}
As for the remaining terms $Q(\gamma,h_\pi)$ and $T_2(\gamma,h_\pi)$, we will identify them essentially with the left-hand side of \eqref{eq:el-est} and the $\Gamma$-term on the right-hand side of \eqref{eq:el-est}, respectively. 
A technical detail one has to take into account here, is that the domain of \eqref{eq:ourELeqn} is the torus $\R / \Z$, whereas the respective domain in \eqref{eq:el-est} is the real line $\R$.
To estimate the other terms, let us introduce for 
$f_1, f_2 \in H^{3/2}_{loc}(\mathbb R, \mathbb R^n)$
the operators
\begin{align*}
 \tilde{Q}_\varepsilon (f_1,f_2)& :=  
 \int_{0}^{1}\int_{\tilde I_\varepsilon}
  \left(\left\langle f_1'(u),\ f_2'(u) \right\rangle -
  \frac { \left\langle f_1(u+w)-f(u), f_2(u+w) -f_2(u) 
  \right\rangle}{w^2} \right) \frac {dw}{w^2}\ du\\
& = \int_{0}^{1}\int_{\tilde I_\varepsilon} \int_0^1 \int_0^1
  \frac {\left( \left\langle f_1'(u),\ f_2'(u) \right\rangle- \left\langle f'_1(u+s_1w), f'_2(u+s_2w)
  \right\rangle	\right) }{w^2} ds_1 ds_2 dw\ du,
\end{align*}
where $\tilde I_\varepsilon := [-1/4, 1/4]\setminus [-\varepsilon, \varepsilon]$,
\[ \tilde Q := \lim_{\eps\searrow0} \tilde Q_\eps, \]
and
\begin{equation}
\begin{aligned}
 &\tilde {T}_2(f_1,f_2) \\ &:=  \int_{0}^{1}\int_{-1/4}^{1/4}
  \left\langle f_1(u+w)-f_1(u),\ f_2(u+w)-f_2(u) \right\rangle \left(\frac 1  {|\gamma(u+w)-\gamma(u)|^4}-
 \frac1{|w|^4}\ \right) dw\ du \\
&= \int_{0}^{1}\int_{-1/4}^{1/4} \int_0^1 \int_0^1 
  \left\langle f_1'(u+s_1 w),\ f_2'(u+s_2 w) \right\rangle w^2\left( \frac 1 {|\gamma(u+w)-\gamma(u)|^4}
- \frac1{|w|^4}\ \right) ds_1\ ds_2\ dw\ du.
\end{aligned}
\end{equation}

Recall that $[-\frac{1}{2},\frac{1}{2}]\backslash (-\varepsilon,\varepsilon)$ is used for the definition of $Q$ %, $T_1$,$T_2$,
and $\tilde I_\varepsilon := [-1/4, 1/4]\setminus [-\varepsilon, \varepsilon]$ is used for $\tilde{Q}$, %$\tilde{T}_1$,$\tilde{T}_2$,
so the difference only contains the set where $\abs{w} > \fracm{4}$, thus $\abs{w}^{-2}$ is not singular. Quantitatively, this reads as
\begin{equation} \label{eq:toRtilde}
 |\tilde Q(\gamma,h) - Q(\gamma,h)| + |\tilde T_2 (\gamma,h) - T_2(\gamma,h)| 
\leq C \|\gamma'\|_{L^2} \|h'\|_{L^2} \leq C \|\gamma'\|_{L^2} \|\phi\|_{L^2},
\end{equation}
where we have used that $\gamma$ is bi-Lipschitz in order to deal with $T_2$.

We now compute
\begin{align}
 \tilde{Q}_\eps(\gamma,h_\pi) &=\int_{0}^{1}\int_{\tilde I_\varepsilon} \int_0^1 \int_0^1
  \br{\left\langle \gamma'(u),\ h_\pi'(u)\right\rangle - \left\langle \gamma'(u+s_1w), h'_\pi(u+s_2w)
  \right\rangle} \d s_1\d s_2\frac{dw}{w^2}\ du \nonumber  \\
 &= \int_{0}^{1}\int_{\tilde I_\varepsilon} \int_0^1 \int_0^1
  \br{\left\langle \gamma'(u),\ \phi_{\pi} \right\rangle - \left\langle \gamma'(u+s_1w),  \phi_\pi(u+s_2w)
  \right\rangle} \d s_1\d s_2\frac{dw}{w^2}\ du \nonumber \\
&\quad - \int_{0}^{1}\int_{\tilde I_\varepsilon} \int_0^1 \int_0^1
  \br{\left\langle \gamma'(u),\ a \right\rangle -\left\langle \gamma'(u+s_1w), a
  \right\rangle} \d s_1\d s_2\frac{dw}{w^2}\ du \label{eq:qtildenullterm} \\
&= \int_{0}^{1}\int_{\tilde I_\varepsilon} \int_0^1 \int_0^1
  \br{\left\langle g'(u),\ \phi(u) \right\rangle - \left\langle g'(u+s_1w), \phi(u+s_2w)
  \right\rangle} \d s_1\d s_2\frac{dw}{w^2}\ du . \label{eq:qtilderhs}
\end{align}
Here we have used that $\int_0^1 \gamma'(u+\tilde w) du =0$ for all $\tilde{w}$, i.e., the term \eqref{eq:qtildenullterm} is constantly zero for any $w$. The term \eqref{eq:qtilderhs} essentially is the $L^2$-pairing of $\lapv g'$ and $\lapv \phi$: More precisely we will show
\begin{equation} \label{eq:toRQ}
 |\tilde{Q}(\gamma,h_\pi) - \int_{\mathbb R} \lapv g' \lapv \phi | \leq C_\gamma \|\phi\|_{L^2}.
\end{equation}
\begin{proof}[Proof of \eqref{eq:toRQ}]
% 
% It was shown by He~\cite{He2000} that for a constant $\alpha$ and any $f_1,f_2 \in C_0^\infty(\R)$,
% \begin{equation*}
% \begin{aligned}
%  &\int_{\mathbb R} \lapv f_1' \lapv f_2' du \\
%  &= \alpha \lim_{\eps\searrow0}\int_{\mathbb R} \int_{\mathbb R\setminus [-\varepsilon, \varepsilon]} \left(\left\langle f_1'(u), f_2'(u) \right\rangle
%  - \frac {\left\langle f_1(u+w)-f(u), f_2(u+2)-f_2(u) \right\rangle} {w^2} \right) \frac {dw}{|w|^2} 
%   du
% \end{aligned}
% \end{equation*}
Using~\eqref{eq:FormulaForNorm}, we have for any $f_1,f_2 \in C_0^\infty(\R)$ as $\eps \to 0$
\begin{equation}\label{eq:Heidentity}
\begin{aligned}
  &\int_{\mathbb R} \lapv f_1' \lapv f_2' du + o(1) \\
 &= c\int_{\mathbb R} \int_{\mathbb R \setminus [-\varepsilon, \varepsilon]} \int_0^1 \int_0^1
  \left(\left\langle f_1'(u), f_2'(u) \right\rangle
 - \left\langle f'_1(u+s_1w), f'_2(u+s_2w) \right\rangle  \right) 
\d s_1\d s_2\frac {dw}{|w|^2} du.
\end{aligned}
\end{equation}
We now connect and \eqref{eq:Heidentity} and \eqref{eq:qtilderhs}. The technical problem is, that the integral of $\phi$ is not a feasible test-function for \eqref{eq:Heidentity}. Let therefore for $\eta_{[-10,10]} \in C_0^\infty ((-11,11))$ and $\eta_{[-10,10]} \equiv 1$ in $[-10,10]$,
\[
 \psi := \brac{u \mapsto \eta_{[-10,10]} \int_0^u \phi} \in C_0^\infty(\R).
\]
Thus, $\psi$ is a feasible test-function for \eqref{eq:Heidentity}, which $\phi$ is not. Moreover,
\begin{equation}\label{eq:l42:psip}
 \psi' = \eta_{[-10,10]} \phi + \eta'_{[-10,10]} \int_0^u \phi \overset{\supp \phi}{\equiv } \phi + \eta'_{[-10,10]} \int_0^u \phi.
\end{equation}
We thus arrive at
\begin{align*}
 \tilde{Q}_\varepsilon(\gamma,h) \;
 &\refeq{qtilderhs}
 \int_{0}^{1}\int_{\tilde I_\varepsilon} \int_0^1 \int_0^1
  \br{\left\langle g'(u),\ \phi(u) \right\rangle - \left\langle g'(u+s_1w), \phi(u+s_2w)
  \right\rangle} \d s_1\d s_2\frac{dw}{w^2}\ du\\
 &\stackrl{$\supp \phi$, $\abs{w}$,\eqref{eq:l42:psip}}\qquad
 \int_\R \int_{\tilde I_\varepsilon} \int_0^1 \int_0^1
  \br{\left\langle g'(u),\ \psi'(u) \right\rangle - \left\langle g'(u+s_1w), \psi'(u+s_2w)
  \right\rangle} \d s_1\d s_2\frac{dw}{w^2}\ du\\ 
 &= \int_\R \int_{\R \backslash [-\varepsilon,\varepsilon]} \int_0^1 \int_0^1
  \br{\left\langle g'(u),\ \psi'(u) \right\rangle - \left\langle g'(u+s_1w), \psi'(u+s_2w)
  \right\rangle} \d s_1\d s_2\frac{dw}{w^2}\ du\\
&\qquad{} - \int_\R \int_{\abs{w} > \fracm{4}} \int_0^1 \int_0^1
  \br{\left\langle g'(u),\ \psi'(u) \right\rangle - \left\langle g'(u+s_1w), \psi'(u+s_2w)
  \right\rangle} \d s_1\d s_2\frac{dw}{w^2}\ du\\
%
  &\refeq{Heidentity}
  \int \langle \lapv g', \lapv \psi' \rangle\\
&\qquad{} - \int_\R \int_{\abs{w} > \fracm{4}} \int_0^1 \int_0^1
  \br{\left\langle g'(u),\ \psi'(u) \right\rangle - \left\langle g'(u+s_1w), \psi'(u+s_2w)
  \right\rangle} \d s_1\d s_2\frac{dw}{w^2}\ du\\
  &\qquad{} + o(1)\\
&\refeq{l42:psip}
\int \langle \lapv g', \lapv \phi \rangle + \int \langle \lapv g', \lapv (\eta'_{[-10,10]} \int_0^{\cdot} \phi) \rangle\\
&\qquad{} - \int_\R \int_{\abs{w} > \fracm{4}} \int_0^1 \int_0^1
  \br{\left\langle g'(u),\ \psi'(u) \right\rangle - \left\langle g'(u+s_1w), \psi'(u+s_2w)
  \right\rangle} \d s_1\d s_2\frac{dw}{w^2}\ du\\
  &\qquad{} + o(1).
\end{align*}
Now, by usual interpolation and/or imbedding of Sobolev spaces, see, e.g. \cite{Tartar2007,Schikorra-Doktor},
\[
 \vrac{ \lapv f}_{2,\R} \aleq \vrac{f}_{2,\R} + \vrac{f'}_{2,\R},
\]
we have
\[
 \| \lapv (\eta'_{[-10,10]} \int_0^{\cdot} \phi) \|_{2,\R} \aleq \vrac{\phi}_1 + \vrac{\phi}_2 \overset{\supp \phi}{\aleq} \vrac{\phi}_2.
\]
Moreover,
\[
\begin{ma}
 &&\int_\R \int_{\abs{w} > \fracm{4}} \int_0^1 \int_0^1
  \br{\left\langle g'(u),\ \psi'(u) \right\rangle - \left\langle g'(u+s_1w), \psi'(u+s_2w)
  \right\rangle} \d s_1\d s_2\frac{dw}{w^2}\ du \\
  &\aleq& \int_{\abs{w} > \fracm{4}} \abs{w}^{-2}\ \vrac{g'}_{2} \vrac{\psi'}_{2}\\
  &\overset{\eqref{eq:l42:psip}}{\aleq}& \vrac{g'}_{2,\R} \brac{\vrac{\phi}_{1,\R} + \vrac{\phi}_{2,\R}} \overset{\supp g}{\aleq} C_\gamma\ \vrac{\phi}_{2,\R}.
 \end{ma}
\]
Thus, we have shown that \eqref{eq:toRQ} holds.
\end{proof}


To estimate $\tilde T_2 (\gamma,h)$ we calculate
\begin{align*}
 \tilde T_2(\gamma,h)
 &= \int_{0}^{1}\int_{-1/4}^{1/4} \int_0^1 \int_0^1
  \left\langle \gamma'(u+s_1 w),\  \phi_\pi(u+s_2 w) 
  + a\right\rangle w^2\cdot\\
  &\qquad\qquad\qquad\qquad{}\cdot\left( \frac 1 {|\gamma(u+w)-\gamma(u)|^4}-
 \frac1{|w|^4}\ \right) ds_1\ ds_2\ dw\ du\\
 &= \int_{0}^{1}\int_{-1/4}^{1/4} \int_0^1 \int_0^1
  \left\langle \gamma'(u+s_1 w),\ \phi(u+s_2 w)- \phi(u+s_1 w) 
+ a \right\rangle w^2 \cdot\\
&\qquad\qquad\qquad\qquad{}\left( \frac 1 {|\gamma(u+w)-\gamma(u)|^4}-
 \frac1{|w|^4}\ \right) ds_1\ ds_2\ dw\ du
\end{align*}
and using
\begin{align*}
 w^2 \left( \frac 1 {|\gamma(u+w)-\gamma(u)|^4}- \frac 1 {|w|^4} \right) 
  & = \frac {w^4} {|\gamma(u+w)-\gamma(u)|^4} 
      \left( \frac {1 - \frac {|\gamma(u+w)- \gamma(u)|^4} {w^4}}{w^2}\right) \\
 & \leq C \frac {1 - \frac {|\gamma(u+w)- \gamma(u)|^2} {w^2}}{w^2} \\
 &  = \frac C 2 \int_{0}^1 \int_0^1 \frac {|\gamma'(u+s_3 w) - \gamma'(u+s_4 w)|^2}{w^2}ds_3 \ ds_4
\end{align*}
we get
\begin{equation} \label{eq:toRT2}
\begin{aligned}
 &|\tilde T_2(\gamma,h)|  \leq \abs a E^{(2)} (\gamma) \\ &\quad  + \int_{0}^{1}\int_{-1/4}^{1/4} \int_{(0,1)^4}
 \frac{ |\gamma'(u+s_1 w)-\gamma'(u+s_2w)| \ |\gamma'(u+s_3 w) - \gamma'(u+s_4 w)|^2 \ 
|\phi(u+s_2 w)|}{w^2}  dsdw\ du \\
& \leq C \|\phi\|_{L^2} + \int_{0}^{1}\int_{-1/4}^{1/4} \int_{(-1,1)^3}
 \frac{ |\gamma'(u+s_1 w)-\gamma'(u)| \ |\gamma'(u+s_3 w) - \gamma'(u+s_4 w)|^2 \ 
|\phi(u)|}{w^2}  dsdw\ du.
\end{aligned}
\end{equation}

From \eqref{eq:toRtilde},\eqref{eq:toRT1},\eqref{eq:toRQ}, and \eqref{eq:toRT2} 
one gets the claim, since $\gamma' = g'$ on $[-1/4,5/4]$.
